% Chapter 10: Discussion

\section{What It Means to Measure the Shape of Science}

The thesis shifts the object of measurement from the content of a field to the geometry of its embedded paper distribution. A field is not treated as a list of topics, a citation community, a productivity unit, or a point in a predefined taxonomy. It is treated as a distribution of title--abstract records in a shared SPECTER2 embedding space. Its morphology is the way that distribution occupies the space: how dispersed it is, how tightly its local neighborhoods are packed, how concentrated those neighborhoods are around hubs, how many spectral directions summarize its variation, and how its position changes over time.

This move makes the analysis complementary to standard science maps. Topic models ask what texts are about. Citation networks ask how works, authors, venues, or fields are connected. Bibliometric indicators often ask how much is produced, cited, or recognized. Embedding-space morphology asks a narrower geometric question: given a representation of scientific documents, what structural shape does a field's collection of papers have?

That question is conditional by design. The measurements are properties of a balanced OpenAlex title--abstract corpus, grouped by OpenAlex primary-topic subfields, embedded with SPECTER2, and measured in the original 768-dimensional representation space. This conditionality is not a weakness to hide; it is what makes the result reproducible. The thesis does not claim to reveal the true map of science. It measures one aspect of scientific organization while keeping the data source, taxonomy, representation model, sampling design, and geometric assumptions visible.

\section{Static Shape, Temporal Change, and Semantic Location}

Chapter~6 establishes the first empirical layer: broad domains have morphological tendencies, but they do not determine morphology. Physical Sciences are generally more dispersed and spectrally extended. Health Sciences are more locally dense and hub-concentrated. Social Sciences and Life Sciences occupy less extreme average positions on several metrics. Yet those averages are blunt summaries. Field-level and subfield-level profiles vary substantially within each domain, and the most atypical subfields are extreme for different reasons: dispersion, novelty, local density, hubness, spectral structure, or combinations of them.

Chapter~7 adds time. Across the population of subfields, median centroid distance, median kNN distance, and PCA \(D_{80}\) decline from the first to the last window, while local distance variability can rise and hubness does not move uniformly. Temporal evolution is therefore not a simple story of fields becoming smaller, denser, or more homogeneous. A subfield may become locally compact while becoming more uneven; another may reduce spectral dimensionality while losing hub concentration; another may move far in centroid space without sharing the same structural trajectory.

Chapter~8 turns morphology into a relation. Field-level profile distances show that taxonomy captures part of the structure, but many nearest morphological neighbors cross domain boundaries. The average field-pair distance increases from 2000--2004 to 2020--2024, so profile space becomes more separated rather than globally convergent. That average divergence coexists with selective convergence: some field pairs move closer through shared changes in local distance and spectral dimensionality, while many others move apart.

Chapter~9 compresses these findings into typologies. Static morphology is weakly clustered. The selected Ward \(k=5\) typology has silhouette \(=0.133\), which supports descriptive anchors rather than natural categories. Temporal trajectories separate more clearly, with silhouette \(=0.396\) and mean subsample ARI \(=0.725\). Subfields are easier to group by how their morphology changes than by their full-period shape.

The centroid and hybrid sensitivity checks complete the interpretation. Clustering subfield centroid embeddings separates somewhat more clearly and aligns more strongly with OpenAlex domains than morphology clustering does. Hybrid clustering becomes highly stable when semantic location receives substantial weight. These results do not undermine the morphology typology; they show why semantic location and morphology must remain separate. A field's position in embedding space is not the same as the shape of its distribution around that position.

\section{Methodological Contributions}

The first contribution is a reproducible corpus design for embedding-based science mapping. The thesis builds a synchronized OpenAlex title--abstract corpus covering 2000--2024, with annual balancing across subfield-year cells. This strategy reduces the dominance of high-volume subfields and makes temporal and cross-sectional comparisons more defensible, while making the limits of the empirical object explicit.

The second contribution is the disciplined use of SPECTER2 as an analytical geometry. Quantitative metrics are computed in the original embedding space rather than on UMAP, PCA, or other low-dimensional projections. This prevents visual maps from silently becoming measurement spaces and keeps representation, measurement, and communication apart.

The third contribution is the reduced eleven-metric morphology core. The metric set covers global dispersion, local density and hubness, spectral structure, and temporal movement. No single metric is asked to summarize a field. Morphology is treated as a profile, which makes it possible to see why a field can be dispersed but locally regular, compact but hub-concentrated, or spectrally simple while temporally novel.

The fourth contribution is the separation of static morphology, temporal morphology, morphological similarity, and semantic-location clustering. A single two-dimensional display can make position, distance, cluster, and category appear to be the same object. The thesis keeps them apart: static profiles describe full-period shape, temporal trajectories describe windowed change, similarity compares profiles between units, and centroid clustering checks semantic location rather than shape.

The fifth contribution is interpretive conservatism in clustering. Static clusters are treated as typological anchors because they are weak. Temporal clusters are treated as dynamic tendencies because they are clearer but still exploratory. Centroid and hybrid clusters are used as sensitivity evidence, not as replacements for morphology. The point is not merely technical; it is a way to use clustering without pretending that every partition is a discovery of natural kinds.

\section{Limitations and Threats to Interpretation}

The main data limitations follow from the corpus design. OpenAlex is open and reproducible, but it is not a neutral census of science. Coverage differs across fields, years, languages, document types, and metadata practices. Primary-topic assignment forces each work into one subfield, even when papers have multiple intellectual affiliations. The English-only filter, article and preprint restriction, title--abstract requirement, annual cap, and sparse subfield-year cells make the corpus a controlled analytical sample rather than the observed universe of publication.

The representation also imposes limits. SPECTER2 embeddings are learned document representations, not ground truth. They compress titles and abstracts into vectors shaped by training data, citation-informed supervision, and scientific-language coverage. They may encode patterns useful for scientific document retrieval while still missing method, evidence, theoretical stance, institutional context, or domain-specific nuance.

The geometric summaries carry their own risks. High-dimensional spaces are affected by anisotropy, distance concentration, local density variation, and hubness. The metrics are summaries of distributions, not full distributional comparisons. Centroid dispersion, kNN distances, hub Gini, and spectral proxies each compress many paper-level relations into a small number of statistics. PCA \(D_{80}\) is an operational proxy for spectral complexity, not a claim about true intrinsic dimension. PCA, PCoA, and UMAP projections are useful for communication, but they are not the spaces in which the quantitative claims are computed.

Clustering is the final limitation to keep in view. The static typology is weak by standard separation diagnostics and should not be treated as evidence that science has five morphological categories. The temporal typology is clearer and more stable, but it remains a descriptive partition of slope profiles. Centroid and hybrid clustering show that semantic location can reconstruct more domain-like structure, which is precisely why those results cannot be interpreted as morphology alone. The validity of the thesis depends on the usefulness of the measured distinctions and on the honesty of their scope, not on any cluster being real in an ontological sense.

\section{Implications and Future Research}

The framework suggests a broader implication for science mapping: fields can be compared by structural shape, not only by topical proximity or citation connectivity. A field can be monitored for broadening, contraction, local density change, hub concentration, or spectral simplification without treating those movements as quality signals. The dynamic typologies in Chapter~9 are especially useful in this sense because they name patterns of movement, compacting, broadening, smoothing, and dimensionalizing, that could guide closer historical or domain-specific investigation.

Any practical use should remain cautious. Morphology should not be used to rank fields, allocate resources mechanically, or infer intellectual maturity. A compact field is not better than a dispersed one. A field with high novelty is not necessarily more innovative. A field whose morphology diverges from others is not necessarily isolated. The value of the framework is diagnostic: it can identify where structural changes or unusual profiles occur, so that substantive experts can ask better questions.

Several extensions follow directly. The same pipeline could be run with other scientific embedding models to test which findings are stable across representations. Multilingual corpora could show how much of the observed morphology is tied to English-language title--abstract records. Full-text or section-aware embeddings could separate abstract-level positioning from methods, results, and discussion structure. Morphology could also be compared with citation networks, journal structures, funding categories, institutional affiliations, or author mobility to ask when shape, topic, and social organization move together.

The point of these extensions is not to build a larger map for its own sake. It is to preserve the key distinction while adding evidence around it. Scientific fields have positions, shapes, trajectories, and relations. Measuring those dimensions separately makes science more legible, and also more resistant to any single, simple map.
